% Generated draft: requires mathematical and source review.
\documentclass[11pt]{article}
\usepackage[margin=25mm]{geometry}
\usepackage{fontspec}
\setmainfont{Latin Modern Roman}
\usepackage{amsmath,amssymb,mathtools,unicode-math}
\setmathfont{Latin Modern Math}
\usepackage{enumitem,needspace,xurl,hyperref}
\setlength{\parindent}{0pt}
\setlength{\parskip}{5pt}
\setlength{\emergencystretch}{3em}
\newcommand{\N}{\mathbb{N}}
\newcommand{\Z}{\mathbb{Z}}
\newcommand{\Q}{\mathbb{Q}}
\newcommand{\R}{\mathbb{R}}
\newcommand{\C}{\mathbb{C}}
\newcommand{\F}{\mathbb{F}}
\newcommand{\A}{\mathbb{A}}
\newcommand{\PP}{\mathbb{P}}
\newcommand{\E}{\mathbb{E}}
\newcommand{\eps}{\varepsilon}
\newcommand{\dd}{\,\mathrm{d}}
\begin{document}
\section*{UnsolvedMath ID: 20002100}
\textbf{Source problem code:} catalogue entry
\subsection{Definitions and mathematical statement}
Question 1.42 records Peter Ozsvath's request for a spectrum theory, in the spirit of C. Manolescu, for Lagrangian intersections. Here a spectrum means a stable homotopy refinement of Floer theory, not an eigenvalue spectrum. The source does not specify the symplectic/Lagrangian class, brane and tangential data, grading, coefficients, analytic hypotheses or invariance requirements. The aid reports partial exact, tangential, parameterized and monotone-truncated constructions, not one universal theorem covering every Lagrangian pair.

Its concrete calibration is the following cotangent graph case. Let \(Q\) be a closed smooth manifold, let \(f:Q\to\R\) be Morse with a Morse--Smale auxiliary metric, and put
\[
 L_0=Q\subset T^*Q,\qquad
 L_\epsilon=\operatorname{graph}(\epsilon\,\dd f),\qquad\epsilon>0.
\]
These are compact exact Lagrangians and intersect transversely at the critical points of \(f\). For sufficiently small \(\epsilon\) and compatible Floer data, the Morse--Floer comparison identifies the Floer complex with the Morse complex. Use \(\F_2\) coefficients to avoid omitted orientation/branes assumptions and the cohomological grading agreeing with Morse degree.

The aid constructs a graph spectrum by the Morse-CJS realization, using the compactified gradient-flow data and coherent realization data that recover the Morse cell attachments of \(Q\). In this convention its stable type is
\[
 \mathcal X_{\mathrm{graph}}(L_0,L_\epsilon)
 \simeq\Sigma^\infty Q_+,
\]
where \(Q_+=Q\sqcup\{*\}\). Consequently
\[
 \widetilde H^k(\mathcal X_{\mathrm{graph}};\F_2)
 \simeq H^k(Q;\F_2)
 \simeq HF^k(L_0,L_\epsilon;\F_2).
\]
This specifies a Morse-calibrated special-case refinement and its normalization. It does not assert that every Floer moduli problem already carries the necessary compactification, framing, twisting and continuation data, nor that conventions using duals or suspension shifts can be compared without adjustment.

For \(Q=\C\mathrm P^2\), a perfect Morse function has generators in degrees \(0,2,4\) and zero differential. Thus the naive wedge of sphere spectra
\[
 \mathbb S^0\vee\mathbb S^2\vee\mathbb S^4
\]
has the same additive \(\F_2\) cochain data as the graph Floer complex. It is nevertheless not stably equivalent to \(\Sigma^\infty(\C\mathrm P^2)_+\). Indeed for the degree-two generator \(x\),
\[
 H^*(\C\mathrm P^2;\F_2)=\F_2[x]/(x^3),\quad|x|=2,
 \qquad\operatorname{Sq}^2(x)=x^2\ne0,
\]
whereas the corresponding operation from degree two to four on the wedge is zero. Stable equivalences preserve Steenrod operations. Therefore matching Floer generators, gradings and differential is insufficient: a proposal compatible with this Morse-CJS calibration must also retain the nontrivial stable attaching information detected by \(\operatorname{Sq}^2\).

The aid gives this as a special-case construction and falsification test. It does not supply a complete universal spectrum construction, a comparison among all existing approaches, or full hypotheses for the exact/tangential/parameterized/monotone-truncated results it mentions.
\subsection{Short English statement}
The small cotangent graph pair has a Morse-calibrated spectrum \(\Sigma^\infty Q_+\) refining its Floer cohomology. For \(Q=\C\mathrm P^2\), a nonzero Steenrod square rules out the naive wedge of generator spheres, even though the additive complex agrees. The general spectrum request retains additional geometric and analytic conditions.
\subsection{Sources}
\begin{itemize}
\item Permanent UnsolvedMath ID \(20002100\), exact public code catalogue entry; local registry \texttt{aim-geometry-notes.json}, \texttt{notes[437]}, canonical tag question.
\item AIM workshop \emph{Holomorphic curves in contact geometry}, Question 1.42, attribution Peter Ozsvath and comparison to C. Manolescu: \url{https://aimath.org/WWN/contactgeom2/contactgeom2.pdf}. No PDF was accessed.
\item The aid invokes Morse-CJS realization and Morse--Floer comparison, and reports several partial spectrum approaches without naming full papers. No original literature entry, proposer/year metadata or authoritative full citation is supplied.
\item Research attempt 1, ``A cotangent-graph calibration for Lagrangian Floer spectra,'' labelled partial\_result, partially\_solved, proved and low novelty confidence. Its summary describes the request as from 2003; the exact-statement verification field merely introduces the canonical record.
\end{itemize}
\end{document}
